\documentclass[12pt,a4paper]{article}
\usepackage[T1]{fontenc}

\title{Introduction to \LaTeX}
\author{Arka Nath\footnote{KUET} \and Dr.\ Sheikh Hasina}

\begin{document}
	
	\maketitle
	\tableofcontents
	
	\section{Introduction}
	
	Khulna University of Engineering \& Technology (KUET), formerly BIT Khulna, 
	is a public technological university located in Khulna, Bangladesh. 
	It emphasizes education and research in engineering and technology. 
	It was founded in 1967 as an engineering college before gradually converting into a university.\\
	
	The academic division of KUET is organized into 20 departments under 4 faculties.\\
	Sixteen departments offer undergraduate courses. Among them, the Department of 
	Industrial Engineering and Management (IEM) provides an Industrial and 
	Production Engineering (IPE) degree to undergraduate students.
	
	\subsection{Departments}
	
	Khulna University of Engineering \& Technology has a background of historical significance. 
	Although this institution has completed an era as a university, its foundation was laid 56 years ago. 
	The university was established in 1967 as ``Khulna Engineering College'' under the University of Rajshahi. 
	During the Liberation War, development activities of the erstwhile Engineering College were suspended. 
	After the independence of Bangladesh in 1971, development works resumed in full swing. 
	Overcoming all constraints, the academic activities continued.
	
	\subsection{Architecture}
	
	After the independence of Bangladesh in 1971, development works resumed in full swing. 
	The academic activities of the institute officially started on June 03, 1974. 
	Despite some limitations, the college was affiliated as a Faculty of Engineering under Rajshahi University. 
	The academic and administrative activities were controlled by the university authority 
	and the Ministry of Education, respectively. 
	The physical infrastructure was constructed by the Public Works Department, while teacher recruitment was accomplished through the Public Service Commission.
	
	\begin{center}
		\begin{tabular}{ c c c }
			cell1 & cell2 & cell3 \\
			cell4 & cell5 & cell6 \\
			cell7 & cell8 & cell9 \\
		\end{tabular}
	\end{center}
	
\end{document}